% !TeX program = XeLaTeX
% !TeX root = ../pujavidhanam.tex
\chapt{श्री-अश्वत्थनारायण-पूजा (अमावास्या-सोमवार-व्रतम्)}

\centerline{\small{(मूलम्—श्रीव्रतरत्नाकरः, प्रथमो भागः)}}

\twolineshloka*
{अमावास्या यदा पार्थ सोमवारान्विता भवेत्}
{तदा पुण्यतमः कालो देवानामपि दुर्लभः}

\input{purvanga/vighneshwara-puja}

\sect{प्रधान-पूजा — श्री-अश्वत्थनारायण-पूजा}

\twolineshloka*
{शुक्लाम्बरधरं विष्णुं शशिवर्णं चतुर्भुजम्}
{प्रसन्नवदनं ध्यायेत् सर्वविघ्नोपशान्तये}

\dnsub{सङ्कल्पः}

ममोपात्त-समस्त-दुरित-क्षयद्वारा श्री-परमेश्वर-प्रीत्यर्थं शुभे शोभने मुहूर्ते अद्य ब्रह्मणः
द्वितीयपरार्धे श्वेतवराहकल्पे वैवस्वतमन्वन्तरे अष्टाविंशतितमे कलियुगे प्रथमे पादे
जम्बूद्वीपे भारतवर्षे भरतखण्डे मेरोः दक्षिणे पार्श्वे शकाब्दे अस्मिन् वर्तमाने व्यावहारिकाणां
प्रभवादीनां षष्ट्याः संवत्सराणां मध्ये \blank\see{app:samvatsara_names} नाम संवत्सरे
\blank\see{app:masa_names}-मासे \blank{}-पक्षे अमावास्यायां सोमवासरयुक्तायां
\blank\see{app:nakshatra_names} नक्षत्र \blank\see{app:yoga_names}-योग \blank\see{app:karanam_names}-करण-युक्तायां
च एवं गुण विशेषण विशिष्टायाम् अस्याम् अमावास्यायां शुभतिथौ
अस्माकं सहकुटुम्बानाम् अवैधव्य-सन्तति-चिरञ्जीविका-सम्पत्-सौभाग्य-सिद्ध्यर्थं
पुत्रपौत्राभि\-वृद्ध्यर्थं क्षेमस्थैर्य-धैर्य-वीर्य-विजय-आयुरारोग्य-ऐश्वर्याभिवृद्ध्यर्थम्
धर्मार्थकाममोक्ष\-चतुर्विधफलपुरुषार्थसिद्ध्यर्थं इष्टकाम्यार्थसिद्ध्यर्थं
मम इहजन्मनि पूर्वजन्मनि जन्मान्तरे च सम्पादितानां ज्ञानाज्ञानकृतमहा\-पातकचतुष्टय-व्यतिरिक्तानां
रहस्यकृतानां प्रकाशकृतानां सर्वेषां पापानां सद्य अपनोदनद्वारा सकल-पापक्षयार्थं
अश्वत्थमूले ब्रह्म-विष्णु-शिव-स्वरूपं श्री-अश्वत्थनारायणम् उद्दिश्य तत्प्रीत्यर्थं
श्री-अश्वत्थनारायण-प्रसादसिद्ध्यर्थं यथाशक्ति-ध्यानावाहनादिषोडशोपचारैः
श्री-अश्वत्थनारायण-पूजां करिष्ये। तदङ्गं कलशपूजां च करिष्ये।

श्रीविघ्नेश्वराय नमः यथास्थानं प्रतिष्ठापयामि। शोभनार्थे क्षेमाय पुनरागमनाय च।\\
(गणपति-प्रसादं शिरसा गृहीत्वा)

\input{purvanga/aasana-puja}

\input{purvanga/ghanta-puja}

\begingroup
\renewcommand{\devaName}{देव}
\input{purvanga/kalasha-puja}
\endgroup

\input{purvanga/aatma-puja}

\input{purvanga/pitha-puja}

\input{purvanga/guru-dhyanam}

\dnsub{अष्टोत्तरशत-अश्वत्थ-प्रदक्षिण-सङ्कल्पः}

अनन्तरं फलैः पुष्पैः भक्ष्यैः वस्त्रादिभिर्वा अश्वत्थवृक्षस्य मूलतो ब्रह्मरूपाय
मध्यतो विष्णुरूपिणे अग्रतः शिवरूपाय वृक्षराजाय ते नमः इति सङ्कल्प्य
अष्टोत्तरशतसंख्याकानि अश्वत्थ-प्रदक्षिणानि करिष्ये।

\sect{षोडशोपचार-पूजा}
\renewcommand{\devAya}{श्री-अश्वत्थनारायणस्वामिने नमः,}
\begin{center}

\dnsub{ध्यानम्}

अतो देवा अवन्तु नो यतो विष्णुर्विचक्रमे पृथिव्यास्सप्तधामभिः, इति मन्त्रेण अश्वत्थमूले विष्णुं ध्यात्वा स्पृष्ट्वा पूजयेत्।

अश्वत्थे वो निषदनं पर्णे वो वसतिः कृता। गोभाज इत्किलासथ यत्सनवद् पुरुषम्, इति मन्त्रेण,

\twolineshloka*
{मूलतो ब्रह्मरूपाय मध्यतो विष्णुरूपिणे}
{अग्रतः शिवरूपाय वृक्षराजाय ते नमः}
इति पुराणमन्त्रेण अश्वत्थस्य षोडशोपचारान् कुर्यात्।

\twolineshloka*
{सर्वव्यापक विश्वेश कृपया दीनवत्सल}
{मयोपपादितां पूजां गृहाणेमां जगत्पते}

\twolineshloka*
{मूलानि ब्राह्मणा यस्य ऋग्वेदोऽङ्गानि पल्लवाः}
{फलानि यज्ञाश्शाखाश्च नारायण दयां कुरु}
\textbf{अश्वत्थमूले ब्रह्मविष्णुशिवस्वरूपं श्री-अश्वत्थनारायणं ध्यायामि, आवाहयामि।}
\medskip

\twolineshloka*
{सूर्यकोटिप्रतीकाश सर्वव्यापिन् जगत्पते}
{आसनं कल्पितं शक्त्या गृहाण परमेश्वर}
\textbf{\devAya{} आसनं समर्पयामि।}
\medskip

\twolineshloka*
{महादेव जगद्व्यापिन् जगदानन्दकारक}
{विष्णुक्रान्तादिसंयुक्तं तव पाद्यं मयार्पितम्}
\textbf{\devAya{} पाद्यं समर्पयामि।}
\medskip

\twolineshloka*
{फलगन्धाक्षतैर्युक्तं पुष्पदर्भासमन्वितम्}
{अर्घ्यं गृहाण भगवन् विष्णो सर्वफलप्रद}
\textbf{\devAya{} अर्घ्यं समर्पयामि।}
\medskip

\twolineshloka*
{कर्पूरैलादिसुरभि शीतलं सलिलं प्रभो}
{उमारमण देवेदम् आचम्यं प्रतिगृह्यताम्}
\textbf{\devAya{} आचमनीयं समर्पयामि।}
\medskip

\twolineshloka*
{पञ्चामृतं मया नीतं पयो दधि घृतं मधु}
{शर्करागुडसंयुक्तं स्नानार्थं प्रतिगृह्यताम्}
\textbf{\devAya{} पञ्चामृतस्नानं समर्पयामि।}
\medskip

\twolineshloka*
{गङ्गा कृष्णा गौतमी च कावेरी च पिनाकिनी}
{सलिलं ताभ्य आनीतं स्नानार्थं प्रतिगृह्यताम्}
\textbf{\devAya{} स्नानं समर्पयामि। स्नानानन्तरम् आचमनीयं समर्पयामि।}
\medskip

\twolineshloka*
{पीतवासस्तव विभो व्याघ्राजिनं मयाऽऽहृतम्}
{वासः प्रीत्या गृहाणेदं पुरुषोत्तम शङ्कर}
\textbf{\devAya{} वस्त्रं समर्पयामि।}
\medskip

\twolineshloka*
{उपवीतं सोत्तरीयं मया दत्तं सुशोभनम्}
{विश्वमूर्ते गृहाणेदं नारायण शिव प्रभो}
\textbf{\devAya{} यज्ञोपवीतं समर्पयामि।}
\medskip

\twolineshloka*
{अश्वत्थ हुतभुग्वास गोविन्दस्य सदाप्रिय}
{अशेषं हर मे पापं वृक्षराज नमोऽस्तु ते}
\textbf{\devAya{} आभरणानि समर्पयामि।}
\medskip

\twolineshloka*
{मलयाचलसम्भूतं घनसारं मनोहरम्}
{विलेपनं सुरश्रेष्ठ चन्दनं प्रतिगृह्यताम्}
\textbf{\devAya{} गन्धान् धारयामि।}
\medskip

\twolineshloka*
{अक्षताश्च सुरश्रेष्ठ कुङ्कुमाक्तास्सुशोभनाः}
{मया निवेदिताः भक्त्या त्वदर्थं परमेश्वर}
\textbf{\devAya{} अक्षतान् समर्पयामि।}
\medskip

\twolineshloka*
{तुलसीजाति-कमल-मल्लिका-चम्पकानि च}
{बिल्वानि हर गोविन्द गृहाण परमेश्वर}
\textbf{\devAya{} पुष्पाणि समर्पयामि।}

\end{center}

\dnsub{चतुस्त्रिंशन्नाम-पर्ण-पूजा}
\begin{multicols}{2}
\begin{enumerate}
  \item ॐ परमात्मने नमः
  \item छन्दःपर्णाय नमः
  \item परमेश्वराय नमः
  \item परब्रह्मणे नमः
  \item नारायणाय नमः
  \item विरिञ्चये नमः
  \item शिवाय नमः
  \item रुद्राय नमः
  \item विभूतये नमः
  \item वनस्पतये नमः
  \item वृक्षराजाय नमः
  \item वेदात्मने नमः
  \item प्रकृतये नमः
  \item देवात्मने नमः
  \item त्रिमूर्त्यात्मने नमः
  \item ऊर्ध्वमूलाय नमः
  \item अव्ययाय नमः
  \item पापघ्नाय नमः
  \item ज्ञानवृक्षाय नमः
  \item भूतावासाय नमः
  \item पुण्याय नमः
  \item शिवङ्कराय नमः
  \item महामूलाय नमः
  \item ऋषिसेविताय नमः
  \item तपःफलाय नमः
  \item अग्निगर्भाय नमः
  \item अरणये नमः
  \item समिदाकराय नमः
  \item सूर्याश्वजाय नमः
  \item शर्मापतये नमः
  \item अधश्शाखाय नमः
  \item अश्वत्थाय नमः
\end{enumerate}
\end{multicols}
\textbf{\devAya{} नानाविध-परिमल-पत्रपुष्पाणि समर्पयामि। एतैः नामभिः केशवादिनामभिश्च पूजां कुर्यात्।}

\begin{center}
\sect{उत्तराङ्ग-पूजा}

\twolineshloka*
{वनस्पतिरसोद्भूतो गन्धाढ्योऽतिमनोहरः}
{आत्रेयस्सर्वदेवानां धूपोऽयं प्रतिगृह्यताम्}
\textbf{\devAya{} धूपमाघ्रापयामि।}
\medskip

\twolineshloka*
{चक्षुर्दः सर्वलोकानां तिमिरस्य निवारणम्}
{आर्तिक्यं कल्पितं भक्त्या गृहाण जगदीश्वर}
\textbf{\devAya{} दीपं दर्शयामि।}
\medskip

\twolineshloka*
{भक्ष्यं भोज्यं लेह्यपेयं चोष्यं खाद्यं मयाऽऽहृतम्}
{प्रीतये परमेशस्य दत्तं मे स्वीकुरु प्रभो}
\textbf{\devAya{} नैवेद्यं निवेदयामि। निवेदनानन्तरम् आचमनीयं समर्पयामि।}
\medskip

\textbf{\devAya{} ताम्बूलं समर्पयामि।}
\medskip

\twolineshloka*
{त्वद्भासा भासते लोकः कोटिसूर्यसमप्रभ}
{नीराजयिष्ये त्वां विष्णो कृपां कुरु जगदीश्वर}
\textbf{\devAya{} कर्पूरनीराजनं दर्शयामि। नीराजनानन्तरम् आचमनीयं समर्पयामि।}
\medskip

\twolineshloka*
{कायेन मनसा वाचा मम जन्मशतार्जितम्}
{पापं प्रशमय स्वामिन् प्रदक्षिणपदे पदे}
\textbf{\devAya{} प्रदक्षिणान् समर्पयामि।}
\medskip

\twolineshloka*
{व्यक्ताव्यक्तस्वरूपाय सृष्टिस्थित्यन्तकारिणे}
{आदिमध्यान्तशून्याय विष्टरश्रवसे नमः}
\textbf{\devAya{} नमस्कारान् समर्पयामि।}
\medskip

\twolineshloka*
{मूलतो ब्रह्मरूपाय मध्यतो विष्णुरूपिणे}
{अग्रतश्शिवरूपाय वृक्षराजाय ते नमः}
\textbf{\devAya{} मन्त्रपुष्पं समर्पयामि। सुवर्णपुष्पं समर्पयामि।}
\medskip

\twolineshloka*
{त्रिमूर्तिरूपिन् विश्वेश वृक्षरूप नमोऽस्तु ते}
{मत्कृतं दुरितं शोध्य ज्ञानं मे देहि शाश्वतम्}
\textbf{प्रार्थनाः समर्पयामि।}
\medskip

\dnsub{अर्घ्यम्}
\resetShloka
ममोपात्त-समस्त-दुरित-क्षयद्वारा श्री-अश्वत्थनारायण-प्रीत्यर्थं
अमासोमवारव्रतोद्यापनान्ते क्षीरार्घ्यप्रदानम् उपायनदानं च करिष्ये।

\twolineshloka*
{लक्ष्मीश सर्वलोकेश सर्वसम्पत्प्रदो भव}
{गृहाणार्घ्यं मया दत्तं कृपया परया मुदा}
\textbf{\devAya{} इदमर्घ्यम् इदमर्घ्यम् इदमर्घ्यम्।}

अनेन अर्घ्यप्रदानेन श्री-अश्वत्थनारायणस्वामी प्रीयताम्।

\dnsub{उपायन-दानम्}

\twolineshloka*
{अमासोमव्रतस्यास्य सम्पूर्णफलहेतवे}
{वायनं द्विजवर्याय सहिरण्यं ददाम्यहम्}

\twolineshloka*
{यन्मया मनसा वाचा नियमात् पूजनं कृतम्}
{सर्वं सम्पूर्णतां यातु तद्विप्राणां प्रसादतः}

\hiranyagarbha

अमासोमवार-पुण्यकाले अस्मिन् मया क्रियमाण-श्री-अश्वत्थनारायण-पूजायां
यद्देयमुपायनदानं तत्प्रत्याम्नायार्थं हिरण्यं श्री-अश्वत्थनारायण-प्रीतिं
कामयमानः मनसोद्दिष्टाय ब्राह्मणाय सम्प्रददे नमः न मम।

अनया पूजया श्री-अश्वत्थनारायणः प्रीयताम्।

इति अमावास्या-सोमवार-व्रतोद्यापनविधिः सम्पूर्णः॥

\end{center}

\sect{अपराध-क्षमापनम्}

\fourlineindentedshloka*
{कायेन वाचा मनसेन्द्रियैर्वा}
{बुद्‌ध्याऽऽत्मना वा प्रकृतेः स्वभावात्}
{करोमि यद्यत् सकलं परस्मै}
{नारायणायेति समर्पयामि}

\centerline{ॐ तत्सद्ब्रह्मार्पणमस्तु।}

\closesub

\ifbool{katha}{\sect{सोमवती-अमावास्या-व्रत-कथा}
\centerline{\small{(मूलम्—श्रीभविष्योत्तरपुराणे, भीष्म-युधिष्ठिर-संवादः)}}

\twolineshloka
{शरतल्पगतं भीष्ममुपगम्य युधिष्ठिरः}
{कृतप्रणामो धर्मात्मा हितं वचनमब्रवीत्} %॥१॥

\uvacha{युधिष्ठिर उवाच}

\twolineshloka
{हतेषु कुरुमुख्येषु भीमसेनेन कोपिना}
{तथाऽपरेषु भूतेषु हतेषु युधि जिष्णुना} %॥२॥

\twolineshloka
{दुर्योधनकुमन्त्रेण जातोऽस्माकं कुलक्षयः}
{न सन्ति भुवि भूपालाः बालवृद्धातुरान् ऋते} %॥३॥

\twolineshloka
{अवशिष्टा वयं पञ्च वंशे भारतसंज्ञके}
{एकातपत्रमपि च राज्यं मह्यं न रोचते} %॥४॥

\twolineshloka
{जीवितेऽपि जुगुप्सा स्यान् मृतेन सुगतिः क्वचित्}
{दृष्ट्वा सन्ततिविच्छेदं सन्तापो हृदयेऽनिशम्} %॥५॥

\twolineshloka
{अश्वत्थाम्ना निर्दग्धस्तूत्तरागर्भसम्भवः}
{अतो मे द्विगुणं दुःखं पिण्डविच्छेददर्शनात्} %॥६॥

\twolineshloka
{किं करोमि क्व गच्छामि पितामह वदाधुना}
{येन सम्पद्यते सद्यस्सन्ततिश्चिरजीविनी} %॥७॥

\uvacha{भीष्म उवाच}

\twolineshloka
{शृणु राजन् प्रवक्ष्यामि व्रतानामुत्तमं व्रतम्}
{यस्यानुष्ठानमात्रेण सन्ततिश्चिरजीविनी} %॥८॥

\twolineshloka
{अमावास्या यदा पार्थ सोमवारयुता भवेत्}
{तस्यामश्वत्थमागत्य पूजयेच्च जनार्दनम्} %॥९॥

\twolineshloka
{अष्टोत्तरशतं कुर्यात्तस्मिन् वृक्षे प्रदक्षिणान्}
{तावत्संख्यान्युपादाय रत्नधातुफलानि च} %॥१०॥

\twolineshloka
{व्रतराजमिदं राजन् विष्णोः प्रीतिकरं शुभम्}
{रुद्रादिलयकर्तारः प्रसीदन्ति तदात्मकाः} %॥११॥

\twolineshloka
{मूर्तित्रयात्मकोऽश्वत्थः सर्वदेवमयस्तथा}
{पूजया वेदवृक्षस्य सर्वेष्टार्थानवाप्नुयात्} %॥१२॥

\twolineshloka
{उत्तरां कारय प्राज्ञ गर्भो जीवितमाप्स्यति}
{भविष्यति गुणैः पुत्रस्त्रिषु लोकेषु विश्रुतः} %॥१३॥

\onelineshloka
{श्रुत्वा पितामहप्रोक्तं प्रत्युवाच युधिष्ठिरः} %॥१४॥

\uvacha{युधिष्ठिर उवाच}

\twolineshloka
{तद्व्रतं व्रतराजाख्यं विस्तरेण प्रकाशय}
{केन प्रकाशितं मर्त्ये केनेदं विहितं विभो} %॥१५॥

\uvacha{भीष्म उवाच}

\twolineshloka
{अस्ति या सर्वविख्याता काञ्ची संज्ञा महापुरी}
{रजताचलसङ्काशसौधसन्धैर्विराजिता} %॥१६॥

\twolineshloka
{सुवेषैर्नागरजनैर्नारीभिरुपसेविता}
{ब्रह्मक्षत्रियविट्शूद्रैः स्वस्वकर्मरतैर्वृता} %॥१७॥

\twolineshloka
{रूपचातुर्ययुक्ताभिर्वेश्याभिः समलङ्कृता}
{अत्र राजा महासेनो बभूवामितविक्रमः} %॥१८॥

\twolineshloka
{तस्मिन् वसति भूपाले जनानन्दकरे शुभे}
{तस्य राज्येऽवसद्विप्रः देवस्वामीति विश्रुतः} %॥१९॥

\twolineshloka
{तस्य स्त्री रूपसम्पन्ना नाम्ना धनवती शुभा}
{यथार्थनामधेया सा लक्ष्मीरिव सुविग्रहा} %॥२०॥

\twolineshloka
{तस्यां सञ्जनयामास सप्त पुत्रान् गुणान्वितान्}
{एकां दुहितरं रम्यां नाम्ना गुणवतीं नृप} %॥२१॥

\twolineshloka
{कृतदाराश्च ते पुत्रा विहरन्ति यथासुखम्}
{कन्या कुमारिका चासीदनुरूपप्रियार्थिना} %॥२२॥

\twolineshloka
{अत्रान्तरे द्विजः कश्चिद्भिक्षार्थं समुपागतः}
{दीप्यमानः स्वतेजोभिर्मूर्तिमानिव पावकः} %॥२३॥

\twolineshloka
{द्वारदेशमुपागत्य प्रयुयोज आशिषं तदा}
{देवस्वामिस्नुषास्सप्त समुत्थाय ससम्भ्रमम्} %॥२४॥

\twolineshloka
{भिक्षां प्रत्येकमानिन्युर्ददुस्तस्मै द्विजन्मने}
{अवैधव्याशिषं प्रादात्ताभ्यः सौभाग्यसम्पदः} %॥२५॥

\twolineshloka
{ततो गुणवती मात्रा प्रहिता सह भिक्षया}
{विप्राय भिक्षां प्रददौ कृत्वा पादाभिवन्दनम्} %॥२६॥

\twolineshloka
{आशिषं च ददौ विप्रः शुभे धर्मवती भव}
{सा वैलक्ष्यं गुणवती श्रुत्वा गत्वा ययौ गृहम्} %॥२७॥

\twolineshloka
{मात्रे निवेदयामास चाशिषं तेन योजिताम्}
{श्रुत्वा धनवती पुत्रीं करे वृत्वा समाययौ} %॥२८॥

\twolineshloka
{प्रणतिं कारयामास पुनस्तस्मै द्विजन्मने}
{तथैवाशिषमुच्चार्य पुनस्तामभ्यनन्दयत्} %॥२९॥

\onelineshloka
{श्रुत्वाऽशिषं धनवती सचिन्ता प्रत्युवाच ह} %॥२९.५॥

\uvacha{धनवती उवाच}

\onelineshloka
{प्रसीद भगवन् विप्र वचनं मेऽवधारय} %॥३०॥

\twolineshloka
{स्नुषाभ्यः प्रणताभ्यो मे त्वया दत्ता वराशिषः}
{अवैधव्यकराः पुत्रसुखसौभाग्यदायकाः} %॥३१॥

\twolineshloka
{सुतायां प्रणतायां मे विपरीतं त्वयोदितम्}
{भद्रे धर्मवती भूया इत्युदीर्य पुनः पुनः} %॥३२॥

\twolineshloka
{आशीः प्रयुक्ता विप्रर्षे कारणं वद सस्वरम्}
{धन्यासि त्वं धनवति प्रख्यातचरिता भुवि} %॥३३॥

\uvacha{द्विज उवाच}

\twolineshloka
{यथायोग्यं प्रयुक्तेयं मयाऽऽशीर्दुहितुस्तव}
{इयं सप्तपदीमध्ये वैधव्यं समवाप्स्यति} %॥३४॥

\twolineshloka
{धर्माचरणमत्यर्थं कर्तव्यमनया शुभे}
{अतो मयाऽऽशीर्दत्तेयं शुभे धर्मवती भव} %॥३५॥

\twolineshloka
{श्रुत्वा धनवती वाक्यं चिन्ताकुलितचेतसा}
{उवाच वचनं दीनं प्रणिपत्य पुनः पुनः} %॥३६॥

\uvacha{धनवती उवाच}

\twolineshloka
{उपायं वेत्सि विप्रेन्द्र वद शीघ्रं दयां कुरु}
{अपायं वेत्सि चेद्विप्र किं करोमि दयानिधे} %॥३७॥

\uvacha{द्विज उवाच}

\twolineshloka
{यदाऽऽगता भवेत् सोमा गृहे वै तव सुन्दरि}
{तस्याः पूजनमात्रेण भवेद्वैधव्यनाशनम्} %॥३८॥

\uvacha{धनवती उवाच}

\twolineshloka
{का सा सोमा त्वया प्रोक्ता का जातिः कुत्र संस्थितिः}
{यस्या आगमनमात्रेण भवेद्वैधव्यभञ्जनम्} %॥३९॥

\onelineshloka
{एतद्वद महाभाग सकलं विस्तरस्व मे} %॥४०॥

\uvacha{द्विज उवाच}

\twolineshloka
{सोमा सा रजकी जात्या स्थितिस्तस्याश्च सिंहले}
{सा चेदायाति ते वेश्म भवेद्वैधव्यभञ्जनम्} %॥४१॥

\twolineshloka
{इत्युक्त्वा ब्राह्मणोऽन्यत्र गतो भिक्षाप्रतीक्षया}
{धनवत्यपि पुत्रेभ्यः प्रोवाच वचनं तदा} %॥४२॥

\twolineshloka
{इयं दुर्ललिता पुत्राः स्वसा गुणवती शुभा}
{सोमागमनमात्रेण भवेद्वैधव्यभञ्जनम्} %॥४३॥

\twolineshloka
{यस्यास्ते भक्तिः पितरि मातुर्वचनगौरवम्}
{स प्रयातु सह स्वस्रा सोमामानयितुं ध्रुवम्} %॥४४॥

\uvacha{पुत्रा ऊचुः}

\twolineshloka
{ज्ञातं मातस्तव स्वान्तं पुत्रीस्नेहवशङ्गतम्}
{यातुं देशान्तरं मातर्दुष्करं प्रख्यापयसि} %॥४५॥

\twolineshloka
{अत्र दुस्तरो हि सिन्धुः शतयोजनविस्तरः}
{अशक्यं गमनं तत्र न क्षमा गमने वयम्} %॥४६॥

\uvacha{देवस्वामी उवाच}

\twolineshloka
{अपुत्रस्सप्तभिः पुत्रैरहं यास्यामि सिंहलम्}
{आनयिष्यामि तां सोमां पुत्रीवैधव्यनाशिनीम्} %॥४७॥

\twolineshloka
{एवंवादिनि सक्रोधे देवस्वामिनि तत्क्षणम्}
{शिवस्वामी कनिष्ठस्तु पुत्रः प्रोवाच सम्मतः} %॥४८॥

\uvacha{शिवस्वामी उवाच}

\twolineshloka
{तात मा वद चैवं त्वं रोषावेशवशं गतः}
{मयि तिष्ठति कः शक्तो गन्तुं द्वीपं हि सिंहलम्} %॥४९॥

\twolineshloka
{इत्युक्त्वा सहसोत्थाय प्रणम्य शिरसा मुदा}
{प्रतस्थे सहितः स्वस्रा द्वीपं सिंहलसंज्ञितम्} %॥५०॥

\twolineshloka
{स कियद्भिर्दिनैर्गत्वा तीरं प्राप सरित्पतेः}
{तर्तुं तमम्बुधिं तत्र प्रयत्नमकरोद्द्विजः} %॥५१॥

\twolineshloka
{स ददर्श सुविस्तीर्णं न्यग्रोधद्रुममन्तिके}
{तत्कोटरे सुखासीना गृध्रराजस्य बालकाः} %॥५२॥

\twolineshloka
{तत्पादपतले स्थित्वा ताभ्यां नीतं तु तद्दिनम्}
{शावकानां कृते गृध्रो गृहीत्वा शिशुभोजनम्} %॥५३॥

\twolineshloka
{प्रादान्न शावका गृध्राद्भोजनं जगृहुर्भृशम्}
{पप्रच्छ बालकान् गृध्रश्चिन्ताकुलितमानसान्} %॥५४॥

\uvacha{गृध्र उवाच}

\twolineshloka
{कथं न भुञ्जते पुत्राः भवन्तः क्षुधिता अपि}
{आनीतं कोमलं मांसं भवद्योग्यमिदं मया} %॥५५॥

\uvacha{बालका ऊचुः}

\twolineshloka
{एतद्वृक्षतले तात मानवावत्र तिष्ठतः}
{अभुञ्जतोस्तयोरद्य कथं भुञ्जामहे वयम्} %॥५६॥

\twolineshloka
{एतच्छ्रुत्वा तु गृध्रस्तु करुणाहृष्टचेतसः}
{तयोरन्तिकमागत्य वचनं समभाषत} %॥५७॥

\uvacha{गृध्र उवाच}

\twolineshloka
{जातस्तु युवयोः कामो वक्तव्योऽत्र मम द्विज}
{क्रियते सर्वथा विप्र भोजनं कर्तुमर्हसि} %॥५८॥

\uvacha{द्विज उवाच}

\twolineshloka
{सिंहलं गन्तुकामोऽहं जलधेः पारमद्य वै}
{सोमानयनमिच्छामि स्वसृवैधव्यनाशनम्} %॥५९॥

\uvacha{गृध्र उवाच}

\twolineshloka
{त्वामहं पारयिष्यामि जलधेः प्रातरेव हि}
{सोमागृहमपि ब्रह्मन् दर्शयिष्यामि सिंहले} %॥६०॥

\twolineshloka
{ततो रात्र्यां व्यतीतायामुदिते च दिवाकरे}
{पारमुत्तारितौ तौ तु गृध्रराजेन वेगिना} %॥६१॥

\twolineshloka
{सिंहलद्वीपमागत्य स्थितौ सोमागृहान्तिके}
{ततः प्रत्यूषसमये संसृज्याङ्कणमण्डलम्} %॥६२॥

\twolineshloka
{लेपनं चक्रतुस्तस्या दिवसे दिवसे शुभम्}
{एवं विदधतोस्तत्र पूर्णस्संवत्सरो गतः} %॥६३॥

\twolineshloka
{स्नुषाः पुत्रान् समाहूय सोमा पप्रच्छ विस्मिता}
{मार्जनं लेपनं कोऽत्र करोति मम कथ्यताम्} %॥६४॥

\twolineshloka
{एकैवाथ जगदुस्सर्वाः कृत्यमिदं हि नः}
{ततः कदाचिद्रजकी निभृता संस्थिता निशि} %॥६५॥

\twolineshloka
{ददर्श ब्राह्मणकन्यां मार्जयन्तीं गृहाङ्गणम्}
{लिम्पन्तमङ्कणं प्राप्तं भ्रातरं च शुचिव्रतम्} %॥६६॥

\twolineshloka
{सोमा पप्रच्छ तौ गत्वा कौ युवां कथ्यतां मम}
{ऊचतुस्तौ तदा सोमामावां ब्राह्मणपुत्रकौ} %॥६७॥

\uvacha{सोमा उवाच}

\twolineshloka
{दग्धास्मि बत नष्टास्मि ब्राह्मणौ गृहमार्जकौ}
{कां गतिं बत यास्यामि पापादस्मादसंशयम्} %॥६८॥

\twolineshloka
{पापजातिरहं ख्याता रजकी सर्वथा द्विज}
{कथं त्वं ब्राह्मणो भूत्वा विरुद्धं मे चिकीर्षसे} %॥६९॥

\uvacha{शिवस्वामी उवाच}

\twolineshloka
{एषा गुणवती नाम्ना स्वसा मम सुमध्यमा}
{अस्याः सप्तपदीमध्ये वैधव्यं सम्प्रपत्स्यते} %॥७०॥

\twolineshloka
{तव सान्निध्यमात्रेण भवेद्वैधव्यभञ्जनम्}
{अतो हेतोः सह स्वस्रा पापकर्म करोमि ते} %॥७१॥

\uvacha{सोमा उवाच}

\onelineshloka
{इतः परं न कर्तव्यम् आयामि तव शासनात्} %॥७२॥

\twolineshloka
{इत्युक्त्वा सा समागत्य स्नुषाभ्यः प्रत्युवाच ह}
{यः कश्चिन्मम गेहेऽस्मिन् म्रियते मानवः क्वचित्} %॥७३॥

\twolineshloka
{तथैव रक्षणीयोऽसौ यावदागम्यते मया}
{कस्यचिद्वचनात्कोऽपि नैव दाह्यः कथञ्चन} %॥७४॥

\twolineshloka
{तथेत्युक्ता स्नुषाभिस्सा सोमा याता महाम्बुधिम्}
{पारमुत्तारयामास क्षणेन द्विजपुत्रकौ} %॥७५॥

\twolineshloka
{स्वयमाकाशमार्गेण सोत्ततार महार्णवम्}
{प्राप्ताः काञ्चीपुरीं सर्वे निमेषात्तत्प्रभावतः} %॥७६॥

\twolineshloka
{सोमां दृष्ट्वा धनवती हृष्टा पूजामथाकरोत्}
{अत्रान्तरे शिवस्वामी तदा देशान्तरात्स्वसुः} %॥७७॥

\twolineshloka
{सदृशं वरमानेतुं जगामोज्जयिनीं प्रति}
{आनिनाय वरं तत्र देवशर्मसुतं सुधीः} %॥७८॥

\twolineshloka
{ब्राह्मणं रुद्रशर्माणं गुणिनं सदृशं स्वसुः}
{ततः सा रजकी सोमा वैवाहिकमथाकरोत्} %॥७९॥

\twolineshloka
{सुलभे च सुनक्षत्रे देवस्वामी तु कन्यकाम्}
{ददौ तस्मै गुणवतीं गुणिने रुद्रशर्मणे} %॥८०॥

\twolineshloka
{ततो वैवाहिकैर्मन्त्रैर्ज्वलमाने हुताशने}
{ततः सप्तपदीमध्ये रुद्रशर्मा मृतस्तदा} %॥८१॥

\twolineshloka
{रुरुदुर्बान्धवाः सर्वे स्थिता सोमा निराकुला}
{आक्रन्दश्च महानासील्लोकानां तत्र पश्यताम्} %॥८२॥

\twolineshloka
{गुणवत्यै च सा तूर्णं व्रतराजप्रभावजम्}
{पुण्यं सङ्कल्प्य विधिवद्ददौ मृत्युनिरासनम्} %॥८३॥

\twolineshloka
{रुद्रशर्मापि तस्माच्च व्रतराजप्रभावतः}
{आससाद तदा जीवं सुप्तवत्सहसोत्थितः} %॥८४॥

\twolineshloka
{एवं विवाहं निर्वर्त्य व्रतराजं निवेद्य च}
{आमन्त्र्य तां धनवतीं सोमा प्रायाद्गृहं प्रति} %॥८५॥

\twolineshloka
{एवं सा रजकी सोमा जीवयित्वा मृतं द्विजम्}
{चचाल हर्षसम्पूर्णा पूर्णकामा स्वमन्दिरम्} %॥८६॥

\twolineshloka
{अत्रान्तरे गृहे तस्याः प्रथमं तनया मृताः}
{पुनस्स्वामी मृतस्तस्या जामाता च ततो मृतः} %॥८७॥

\twolineshloka
{आगच्छन्त्यास्तदा तस्याः सोमवारान्विता तिथिः}
{अमावास्या बभूवाथ मृतसञ्जीविनी तिथिः} %॥८८॥

\twolineshloka
{सा ददर्श जलोपान्ते वृद्धां काञ्चित् स्त्रियं तदा}
{तूलभारभराक्रान्तां क्रन्दमानां सुदुःखिताम्} %॥८९॥

\uvacha{वृद्धा उवाच}

\twolineshloka
{अवतारय मे पुत्रि तूलभारं शिरःस्थितम्}
{एतद्भारभराक्रान्तां क्रन्दमानां सुदुःखिताम्} %॥९०॥

\uvacha{सोमा उवाच}

\twolineshloka
{अमावास्याद्य हे वृद्धे सोमवारयुता तिथिः}
{तूलकं न स्पृशाम्यद्य नियमोऽयं मया कृतः} %॥९१॥

\twolineshloka
{पुनर्ददर्श यान्तीं सा मूलभारवतीं स्त्रियम्}
{सा चोवाच ततः पुत्रि मूलभारो महानिति} %॥९२॥

\uvacha{वृद्धा उवाच}

\onelineshloka
{अवतीर्य क्षणं तिष्ठ सङ्के यास्याम्यहं तव} %॥९३॥

\uvacha{सोमा उवाच}

\onelineshloka
{अद्य मूलं तथा तूलं न स्पृशामि कदाचन} %॥९४॥

\twolineshloka
{ततोऽश्वत्थतरुं प्राप्य नदीतीरभवं शुभम्}
{स्नात्वा देवं समभ्यर्च्य शर्करामिः प्रदक्षिणान्} %॥९५॥

\onelineshloka
{सा चकार महाभागा तथैवाष्टोत्तरं शतम्} %॥९६॥

\uvacha{भीष्म उवाच}

\twolineshloka
{यदा प्रदक्षिणशतं कृतं शर्कराहस्तया}
{तदा जीवितवन्तस्ते जामातृपतिपुत्रकाः} %॥९७॥

\twolineshloka
{नगरं श्रीसमाकीर्णं तद्गृहं च विशेषतः}
{अथ याता महाभागा सोमा स्वभवनं प्रति} %॥९८॥

\twolineshloka
{जीवितं वीक्ष्य भर्तारं पुत्रान् जामातरं तथा}
{अभिज्ञातान् समासाद्य जगाम कृतकृत्यताम्} %॥९९॥

\twolineshloka
{प्रणिपत्य स्नुषाः सर्वाः प्रपच्छुस्तां तपस्विनीम्}
{मृतास्ते तनया देवि पतिजामातृबान्धवाः} %॥१००॥

\onelineshloka
{जीवितास्ते कथं सोमे मृतास्ते च कथं वद} %॥१०१॥

\uvacha{सोमा उवाच}

\twolineshloka
{गुणवत्यै मया दत्तं व्रतराजस्य पुण्यकम्}
{मृतास्ते तद्विपाकेन पतिजामातृपुत्रकाः} %॥१०२॥

\twolineshloka
{तूलकं न मया स्पृष्टं मूलकं च तथा स्नुषाः}
{अश्वत्थे देवमभ्यर्च्य कृतास्तत्र प्रदक्षिणाः} %॥१०३॥

\twolineshloka
{शर्कराहस्तया तत्र कृतमष्टोत्तरं शतम्}
{जीवितास्तत्प्रसादेन हृद्याजामातृपुत्रकाः} %॥१०४॥

\twolineshloka
{सर्वाभिः क्रियतामद्य व्रतराजो विधानतः}
{भविष्यति न वैधव्यं सौभाग्यं सम्भविष्यति} %॥१०५॥

\twolineshloka
{स्नुषास्ताः कारयामास तथा सोमा व्रतेश्वरम्}
{भुक्त्वा भोगान्बहूंस्तत्र पुत्रपौत्रादिभिस्सह} %॥१०६॥

\twolineshloka
{तैश्च सर्वैः परिवृता विष्णुलोकमवाप सा}
{इत्येतत्कथितं पार्थ विस्तरेण मया तव} %॥१०७॥

\uvacha{युधिष्ठिर उवाच}

\twolineshloka
{माहात्म्यं व्रतराजस्य को विधिर्वेद विस्तरात्}
{कर्तव्यं केन भीष्मेदं नारीभिः पुरुषेण वा} %॥१०८॥

\uvacha{भीष्म उवाच}

\twolineshloka
{अमावास्या यदा पार्थ सोमवारान्विता भवेत्}
{तदा पुण्यतमः कालो देवानामपि दुर्लभः} %॥१०९॥

\twolineshloka
{प्रातरुत्थाय विधिना स्नानं कार्यं जलाशये}
{स्नात्वा मौनेन कौशेयं परिधाय व्रती ततः} %॥११०॥

\twolineshloka
{गत्वा चाश्वत्थवृक्षस्य समीपं कुरुनन्दन}
{अश्वत्थमूले कर्तव्यं विष्ण्वाद्यमरपूजनम्} %॥१११॥

\twolineshloka
{व्यक्ताव्यक्तस्वरूपाय सृष्टिस्थित्यन्तकारिणे}
{आदिमध्यान्तहीनाय विष्टरश्रवसे नमः} %॥११२॥

\twolineshloka
{इति विष्णुमुखान्देवानर्चयेद्विधिपूर्वकम्}
{छन्दोमये वृक्षराजे सर्वे देवाः समाहिताः} %॥११३॥

\twolineshloka
{मूर्तित्रयात्मन्यश्वत्थनारायणसमाह्वये}
{पूजिते विश्वरूपेऽसौ पूज्यते भगवान् हरः} %॥११४॥

\twolineshloka
{अश्वत्थे वो निषदनं त्वयारात्तेऽग्निरस्तु ते}
{समगृभ्णीत यदुर्ध्वमूलं तद्विष्णोरिति मन्त्रतः} %॥११५॥

\twolineshloka
{आ त्वा वहन्तु देवा वै ब्रह्मज्ञानमिति ह्यपि}
{पूजयेत्सर्वदेवांश्च विष्णुं सर्वात्मकं तथा} %॥११६॥

\twolineshloka
{मूलतो ब्रह्मरूपाय मध्यतो विष्णुरूपिणे}
{अग्रतः शिवरूपाय वृक्षराजाय ते नमः} %॥११७॥

\twolineshloka
{एवं सम्पूज्य देवेशं पीतवस्त्राक्षतैः फलैः}
{कुसुमैर्विविधैश्चैव भक्ष्यभोज्यैस्तथाविधैः} %॥११८॥

\twolineshloka
{अश्वत्थपूजनं कार्यं प्रोक्तमन्त्रेण पाण्डव}
{अश्वत्थ हुतभुग्वास गोविन्दस्य सदाप्रिय} %॥११९॥

\twolineshloka
{अशेषं हर मे पापं वृक्षराज नमोऽस्तु ते}
{प्राणिसञ्जीवकं स्वामिन् भूतावास वनस्पते} %॥१२०॥

\twolineshloka
{अश्वत्थ रुद्ररूपोऽसि वृक्षरूप नमोऽस्तु ते}
{एवं पूजाविधिं कृत्वा ततः कुर्यात्प्रदक्षिणान्} %॥१२१॥

\twolineshloka
{मौक्तिकैः काञ्चनै रौप्यैर्हारकैर्मणिभिस्तथा}
{कांस्यपात्रैस्तानपात्रैर्भक्ष्यपूर्णैः पृथक्पृथक्} %॥१२२॥

\twolineshloka
{गृहीत्वा भ्रमणं कार्यं यावदष्टोत्तरं शतम्}
{समर्पितं च यद्द्रव्यमर्पयेद्द्विजवे शुभम्} %॥१२३॥

\twolineshloka
{सुवासिन्यश्च सम्पूज्याः सोमासन्तुष्टिहेतवे}
{दद्याच्चान्नं तु विप्रेभ्यः स्वयं भुञ्जीत वाग्यतः} %॥१२४॥

\onelineshloka
{एष ते कथितो भूप व्रतराजविधिर्मया} %॥१२५॥

\twolineshloka
{द्रौपदीं च सुभद्रां च कारयस्व तथोत्तराम्}
{उत्तरागर्भसंस्थस्तु जीवितं च लभेच्चिरात्} %॥१२६॥

\uvacha{युधिष्ठिर उवाच}

\twolineshloka
{या स्वल्पविभवा नारी काञ्चनाद्यैर्विनाकृता}
{सा कथं लभते पूर्णं व्रतराजफलं वद} %॥१२७॥

\uvacha{भीष्म उवाच}

\twolineshloka
{फलैः पुष्पैस्तथा भक्ष्यैर्वस्त्राद्यैरपि पाण्डव}
{कुर्यात्प्रदक्षिणावर्तं सापि पूर्णं लभेत्फलम्} %॥१२८॥

\twolineshloka
{व्रतमिदमखिलं नरेन्द्र विष्णोर्विबुधगणस्य च कीर्तितं मयाद्य}
{पतिसुतधनमिच्छती पुरन्ध्री सपदि करोतु फलानि सन्तु नित्यम्} %॥१२९॥

\dnsub{अथ उद्यापनविधिः}

\twolineshloka
{शृणु पार्थ प्रवक्ष्यामि उद्यापनविधिं शुभम्}
{यं विना पूर्णता न स्याद्व्रतराजस्य वै नृप} %॥१३०॥

\twolineshloka
{कारयेत्सर्वतोभद्रं तन्मध्ये कुम्भमुत्तमम्}
{वृक्षराजं सुवर्णस्य पञ्चरत्नस्य वेदिकाम्} %॥१३१॥

\twolineshloka
{तन्मूले प्रतिमां विष्णोः सौवर्णीं च चतुर्भुजाम्}
{लक्ष्मीवाहनसंयुक्तां माषद्वयफलावधि} %॥१३२॥

\twolineshloka
{ब्रह्मरुद्रादिदेवानां यथाशक्ति तु कारयेत्}
{रिक्तस्य सर्वदेवात्मा विष्णुरेव विशिष्यते} %॥१३३॥

\twolineshloka
{उपचारैरनन्तैश्च यथाविभवविस्तरैः}
{नैवेद्यैः पुष्पधूपैश्च दीपैश्च परितः स्थितः} %॥१३४॥

\twolineshloka
{रात्रौ जागरणं कुर्यात्प्रभाते होममाचरेत्}
{समिद्भिः पैप्पलाभिश्च पायसेन तिलैस्तथा} %॥१३५॥

\twolineshloka
{इदं विष्णविति मन्त्रेण हुत्वा च प्रणवेन च}
{विष्ण्वष्टाक्षरमन्त्रेण शिवपञ्चाक्षरेण वा} %॥१३६॥

\twolineshloka
{सर्वदैवतमन्त्रैश्च विष्णुरत्र विशिष्यते}
{मूर्तित्रयात्मकाश्वत्थवृक्षं पूर्णाहुतिं चरेत्} %॥१३७॥

\twolineshloka
{आचार्यं पूजयेत्पश्चाद्गां च दद्यात्पयस्विनीम्}
{ब्राह्मणं वस्त्रभूषाद्यैस्सदस्यं च प्रपूजयेत्} %॥१३८॥

\twolineshloka
{ब्रह्मणे कलशश्चैको विष्णोरेकः प्रधानकः}
{रुद्रादिसर्वदेवानां दशानुगमतः स्मृताः} %॥१३९॥

\twolineshloka
{एवं द्वादश दद्याद्वै तदशक्तो द्वयं तथा}
{ऋत्विजोऽत्र क्रमात्पूज्या घृतपायसभोजनैः} %॥१४०॥

\twolineshloka
{उपवीतानि वस्त्राणि तेभ्यो दद्याच्च दक्षिणाम्}
{प्रणम्य दण्डवद्भूमौ प्रार्थयित्वा विसर्जयेत्} %॥१४१॥

\twolineshloka
{एवं द्वादशवर्षेषु मध्ये चादौ समाचरेत्}
{कृत्वा ह्युद्यापनं सम्यग्व्रतराजफलं भवेत्} %॥१४२॥

\twolineshloka
{सर्वं निवेदयेत्पीठमाचार्याय सदक्षिणम्}
{अच्छिद्रं वाचयेत्पश्चात्स्वयं भुञ्जीत वाग्यतः} %॥१४३॥

॥ इति श्रीभविष्योत्तरपुराणे भीष्मयुधिष्ठिरसंवादे सोमवती-अमावास्या-व्रतकथा सम्पूर्णा ॥
}{}

\closesection
